\section{点列收敛}
\label{sec:topo-convergence}

实数列$x_n\to x$的定义是：对每个$\varepsilon>0$，存在$N$，使$n\geq N$时$\lvert x_n-x\rvert<\varepsilon$。条件$\lvert x_n-x\rvert<\varepsilon$就是$x_n\in(x-\varepsilon,x+\varepsilon)$，所以这句话等于说：对$x$的邻域基中每个区间，点列最终停在其中。把邻域基换成全部邻域，就得到不需要距离的收敛定义。

\begin{definition}[点列收敛]
\label{def:topo-sequence-convergence}
设$X$是拓扑空间，$(x_n)_{n\geq1}$是$X$中的点列，$x\in X$。若对$x$的每个邻域$N$，都存在$n_0$使$n\geq n_0$时$x_n\in N$，就称$(x_n)$\term{收敛}（converge）到$x$，记作$x_n\to x$，并称$x$为点列的一个\term{极限}（limit）。
\end{definition}

由于$x$的每个邻域都包含邻域基中的成员，只需对某个邻域基的成员检验这个条件。在$\mathbb R$上取邻域基$\{(x-\varepsilon,x+\varepsilon):\varepsilon>0\}$，定义~\ref{def:topo-sequence-convergence}就逐字回到实数列极限的$\varepsilon$-$N$定义，所以通常拓扑下的收敛与卷一的极限完全一致。

\paragraph{拓扑决定哪些点列收敛}
收敛完全由邻域决定，因此随拓扑而变。在离散拓扑中，$\{x\}$是$x$的邻域，$x_n\to x$意味着点列最终恒等于$x$。在平凡拓扑中，$x$的唯一邻域是整个空间，每个点列都收敛到每个点，极限可以不唯一。极限何时唯一，取决于不同点能否用不交的邻域分开，这是第~\ref{sec:topo-separation}小节的问题；在那之前，我们只说“一个极限”，而不说“极限”。

\paragraph{子列与闭集}
实数列的两条基本性质只用到邻域，因此在一般拓扑空间中照样成立：收敛点列的子列收敛到同一个点，闭集不会让其中点列的极限跑到集合外面。

\begin{theorem}[子列与闭包]
\label{thm:topo-subsequence-closure}
设$x_n\to x$。
\begin{enumerate}
  \item 对任意子列$(x_{n_k})$（$n_1<n_2<\cdots$），$x_{n_k}\to x$；
  \item 若所有$x_n$都属于集合$A$，则$x\in\overline A$；特别地，若$A$是闭集，则$x\in A$。
\end{enumerate}
\end{theorem}
\begin{proof}
对$x$的邻域$N$，取$n_0$使$n\geq n_0$时$x_n\in N$。由于$n_k\geq k$，当$k\geq n_0$时$n_k\geq n_0$，于是$x_{n_k}\in N$，第一条成立。

对第二条，$x$的每个邻域$N$都含有某个$x_n$，而$x_n\in A$，所以$N$与$A$相交，$x$是$A$的闭包点。$A$闭时$\overline A=A$。
\end{proof}

\paragraph{点列不能反推闭包}
在$\mathbb R$中，定理~\ref{thm:topo-subsequence-closure}的第二条可以反过来：$\overline A$中的每个点都是$A$中某个点列的极限，这是卷一闭集序列判据的内容。下面的例子说明，在一般拓扑空间中这个逆命题不成立，这也是第~\ref{sec:topo-dense-sets}小节不用点列定义稠密性的原因。

\begin{example}[余可数拓扑]
\label{ex:topo-cocountable}
在$\mathbb R$上规定：$U$为开集当且仅当$U=\varnothing$或$\mathbb R\setminus U$是可数集（有限或可数无限）。可数集的任意交、有限并仍可数，所以这些$U$对任意并和有限交封闭，构成一个拓扑，称为\term{余可数拓扑}（cocountable topology）。它的闭集是$\mathbb R$本身与全部可数集。

先看收敛的点列。设$x_n\to x$，令$C=\{x_n:x_n\neq x\}$，它是可数集，所以$\mathbb R\setminus C$是含$x$的开集。点列最终落在其中，即$n$充分大时$x_n\notin C$，也就是$x_n=x$。所以在这个拓扑中，只有最终恒定的点列才收敛。

再取$A=\mathbb R\setminus\{0\}$。包含$A$的闭集只能是$\mathbb R$，因为可数集不能包含不可数集$A$；于是$\overline A=\mathbb R\ni0$。但$A$中的点列若收敛，最终恒等于某个$A$中的点，极限不可能是$0$。所以$0\in\overline A$，却不是$A$中任何点列的极限。同时，$A$中收敛点列的极限都在$A$内，而$A$不是闭集：“包含点列的所有极限”在这里推不出闭性。
\end{example}

问题出在邻域的数量上。在$\mathbb R$的通常拓扑中，$0$的可数个邻域$(-1/k,1/k)$就能代表所有邻域，从每个中选一点即可造出趋于$0$的点列。余可数拓扑中$0$没有这样的可数邻域基：若$N_1,N_2,\dots$是$0$的邻域，各含开集$U_k\ni0$，则$\bigcup_k(\mathbb R\setminus U_k)$是可数个可数集之并，仍可数，于是可取$y\neq0$属于所有$U_k$；开集$\mathbb R\setminus\{y\}$是$0$的邻域，却不包含任何$N_k$。

第~\ref{sec:topo-countability}小节把“有可数邻域基”定义为第一可数性；在度量空间中这一条件自动成立，闭包与闭集的序列判据在第~\ref{chap:metric-spaces}章证明。
